\section{Implementation and side channels}\label{sec:impl}

Revision~1 described the Rust crate \art{crypto:eigen_seta/} (about 4,300
lines plus tests), and \S\ref{sec:impl}--\S\ref{sec:results} still describe
and measure that historical build unless they say otherwise. The maintained
implementation is now the sibling repository \art{eigenx:}~\cite{EigenX}.
\S\ref{sec:impl-v2} lists what changed. The crate is generic over the limb
count, so one code base serves all parameter sets. Variable-time big integers (\texttt{num-bigint}) appear only
in key generation and in public constants.

\subsection{Arithmetic}
\begin{itemize}
\item \emph{$\F_p$.} Montgomery multiplication. For our primes $p+1\equiv0
  \bmod 2^{64Z}$, where $Z$ is 16, 24 and 32 limbs out of 21, 32 and 42, so the
  reduction skips the zero limbs of $p+1$ (``Montgomery-friendly'' primes) and
  costs $N(N-Z)$ instead of $N^2$ word products.
\item \emph{Inversion.} Bernstein--Yang safegcd~\cite{BY19}, generalising the
  structure of libsecp256k1's \texttt{modinv64} to any limb count, with the
  proven divstep bound. It is 40 to 140 times faster than Fermat inversion
  at L1--L5 (Table~\ref{tab:field}).
\item \emph{$\F_{p^2}$.} Karatsuba multiplication, inversion via the norm,
  and the constant-time square root of Adj and Rodr\'iguez-Henr\'iquez~\cite{AR14}.
\item \emph{Curves and isogenies.} x-only Montgomery ladders, including a
  three-point ladder, with Okeya--Sakurai $y$-recovery~\cite{OS01}; SIDH-style
  3-isogeny chains with optimal strategies~\cite{CLN16}; Weierstrass
  arithmetic with constant-time selection between the doubling and addition
  formulas; and the gluing, Richelot and splitting maps of~\cite{CD23} in
  Mumford coordinates, with batched inversions. The Weil pairing uses a
  projective Miller loop over the public bits of $3^F$. Discrete logarithms in
  $\mu_{3^F}$ are computed digit by digit with constant-time comparisons.
\end{itemize}

\subsection{Constant-time decapsulation}
Decapsulation is split into two functions so that each can be checked
separately.
\begin{itemize}
\item \texttt{decaps\_prepare} depends only on $(\mathit{pk},\mathit{ct})$.
  It parses both, lifts $S$ and $T$, and computes both canonical bases. It
  runs in variable time.
\item \texttt{decaps\_core} contains everything secret: the Kani chain with
  its fixed schedule, the identification over the 8 sign and order choices,
  the pairings and DLPs, the 3-isogeny chain, the re-encryption, and the
  selection.
\end{itemize}
In the core, no branch and no memory index depends on secret data, and all
loop counts are public. Selections use \texttt{subtle}'s constant-time
primitives, and every mask passes through \texttt{core::hint::black\_box}.
Scalars have fixed widths. Secret scalars and secret byte buffers are wiped on
drop. The only loop whose trip count depends on data, the canonical-basis
search, runs on public curves and is capped. A crafted public key or
ciphertext makes it return a refusal rather than spin.

\paragraph{A compiler-introduced leak.} The masks in safegcd's final
normalisation (add $m$ if negative, negate if requested) are computed
arithmetically, as \texttt{r + (m \& mask)}. LLVM compiled the add into a
conditional that skipped the work when the mask was zero, so the instruction
count depended on the sign of a secret-dependent intermediate. That leak
affected every field inversion. The source had no branch. The
instruction-count probe of \S\ref{sec:results} found it, and annotating the
probe at line level located it. Passing the masks through
\texttt{black\_box} fixed it. This is a reminder that constant-time code must
be checked on the binary.

\subsection{Residual issues}
\begin{itemize}
\item Key generation is variable time: rejection sampling, Cornacchia and
  primality tests. It must run where no attacker can time it.
\item Allocation: the core allocates vectors whose sizes are fixed by the
  parameter set. glibc's allocator then executes a number of instructions
  that depends on the heap state left by earlier variable-time work, though
  not on secrets. A production implementation should use fixed buffers.
\item Points and field elements on the stack are not wiped, only scalars and
  byte buffers.
\item The memory-address side of constant time is argued, not measured: there
  are no secret-indexed tables. Power, EM and fault attacks are out of scope.
  The FO check catches faults that corrupt $s$, but not faults that skip the
  check.
\item Constant time is verified for one compiler (rustc~1.97, release
  profile) and one target (x86-64). Other builds need the probe re-run.
\end{itemize}

\subsection{Tests}
The crate's tests compare the field arithmetic with a big-integer reference
at 1, 2, 9, 21 and 42 limbs, the chain with direct endomorphism evaluation,
and the pairing with bilinearity. They also run round trips at Toy16 and
Toy32. Known-answer tests pin SHA3-256 digests of the public key, ciphertext
and shared secret produced from fixed seeds at Toy16, Toy32 and L1. The
robustness tests feed decapsulation 48 adversarial ciphertexts
(\S\ref{sec:cca}) and encapsulation 14 adversarial public keys; every call
terminates without a panic. A
regression test covers a key-generation hang that we fixed: Tonelli--Shanks
accepted $z=0$ as a non-residue.

\subsection{The maintained implementation (revision 2)}\label{sec:impl-v2}
The following changes in \art{eigenx:} matter for this paper.
\begin{itemize}
\item \emph{Theta-coordinate decapsulation.} The $(2^E,2^E)$-chain of
  decapsulation is a level-2 theta chain ported from Theta-Rust
  (Dartois--Maino--Pope--Robert~\cite{DMPR23};
  \art{eigenx:src/theta.rs}, \art{eigenx:src/chain.rs}). The Mumford chain of
  \S\ref{sec:impl} is kept for variable-time chain inspection only. The
  identification of $\gamma$ on $\EC[d]$ now matches $x$-coordinates per image
  and selects in constant time among the $y$-sign patterns.
\item \emph{Parameters.} The Q3 and Q5 sets, the NUMS prime search of (P4),
  and parallel key-generation attempts. The toy sets are no longer shipped.
\item \emph{Cryptanalysis code.} A cost model (\art{eigenx:src/security.rs})
  and runnable attacks (\art{eigenx:src/cryptanalysis/}): key derivation from
  $f$, the near-miss experiment, meet in the middle on $\varphi$,
  mask-guess/Kani scaffolding, and the public-key degree test of
  Lemma~\ref{lem:pktest} (new in this revision).
\item \emph{A field-arithmetic defect, fixed in this revision.} The
  interleaved products \texttt{Fp::mul2} and \texttt{Fp::mul3}, which
  $\F_{p^2}$ multiplication and squaring use, always took the
  Montgomery-friendly reduction. For a modulus with $p+1\not\equiv0\bmod2^{64}$
  they returned wrong products. The test \texttt{f1} of
  \art{eigenx:tests/field.rs}, on the field of $2^{61}-1$, failed on every run
  on the main branch and passes after the fix
  (\art{crypto:research/eigen_seta_assessment_20261017/records/}). Every Eigen-S\'eta prime
  is Montgomery-friendly, so no parameter set was affected. The fix routes
  other moduli to the generic product, as single multiplication already did.
\end{itemize}
Constant-time evidence (\S\ref{sec:results}) was produced on the historical
Mumford build and has not been repeated on the theta build;
\art{eigenx:SIDE_CHANNELS.md} is the author's assessment of the current code.

